\section{Cara Menggunakan Buku Ajar}
\label{sec:cara-menggunakan}

Buku ini dapat dibaca dengan dua cara: berurutan dari awal sampai
akhir, atau per bab sesuai minggu perkuliahan. Kedua cara itu sah.
Yang tidak disarankan adalah membaca bab secara acak, karena bab-bab
materi saling bergantung seperti dijelaskan pada
Bagian~\ref{sec:tujuan-buku}.

\subsection{Struktur Setiap Bab}

Setiap bab materi mengikuti kerangka yang sama. Mengenal kerangka ini
akan mempercepat Anda menemukan bagian yang dibutuhkan.

\begin{table}[H]
    \centering
    \caption{Kerangka baku setiap bab dan cara memanfaatkannya.}
    \label{tab:kerangka-bab}
    \small
    \begin{tabular}{@{}p{3.8cm}p{9.6cm}@{}}
        \toprule
        \textbf{Bagian bab} & \textbf{Cara memanfaatkannya} \\
        \midrule
        Sub-CPMK & Baca lebih dahulu. Inilah yang akan diukur pada
        bab ini. \\
        \addlinespace
        Tujuan Pembelajaran & Jadikan daftar periksa: setelah selesai
        membaca, Anda harus mampu melakukan semua yang tertulis. \\
        \addlinespace
        Materi Pokok (section) & Bacaan utama. Setiap konsep disertai
        contoh dan tabel. \\
        \addlinespace
        Contoh Program & Baca kode, lalu \emph{jalankan sendiri}.
        Jangan hanya membaca. \\
        \addlinespace
        Praktikum & Kerjakan di laboratorium atau di komputer
        sendiri; inilah bagian penilaian kinerja. \\
        \addlinespace
        Latihan & Kerjakan tanpa melihat kunci. Latihan menulis
        program lebih penting daripada latihan membaca. \\
        \addlinespace
        Rangkuman \& Checklist & Periksa diri. Tandai butir yang belum
        dikuasai dan ulangi bagian terkait. \\
        \addlinespace
        Evaluasi & Ukur pencapaian. Kerjakan setelah seluruh bagian
        lain selesai. \\
        \bottomrule
    \end{tabular}
\end{table}

\subsection{Alur Belajar yang Disarankan}

Untuk setiap pertemuan, ikuti alur berikut.

\begin{enumerate}
    \item \textbf{Sebelum kelas (30--45 menit).} Baca bagian
          Sub-CPMK, Tujuan Pembelajaran, dan materi pokok bab yang
          bersangkutan. Catat istilah yang belum dipahami.
    \item \textbf{Saat kelas.} Fokus pada praktik dan tanya jawab.
          Karena teori sudah dibaca di rumah, waktu kelas digunakan
          untuk menulis dan menjalankan program.
    \item \textbf{Setelah kelas (60--90 menit).} Kerjakan contoh
          program, praktikum, dan latihan. Tuliskan hasilnya pada
          jurnal praktikum.
    \item \textbf{Menjelang asesmen.} Kerjakan bagian Evaluasi dan
          periksa Checklist Pencapaian Kompetensi.
\end{enumerate}

\begin{aktivitasbox}
Aturan praktis yang perlu dipegang: \textbf{jangan pernah membaca
kode assembly tanpa menjalankannya.} Pemahaman assembly tumbuh dari
mengamati perubahan register dan flag pada setiap langkah eksekusi.
Karena itu, siapkan emulator sejak hari pertama.
\end{aktivitasbox}

\subsection{Portofolio dan Jurnal Praktikum}

Sebagian besar penilaian mata kuliah ini bersifat kinerja, bukan
hafalan. Karena itu, Anda perlu memelihara dua dokumen sepanjang
semester.

\begin{description}
    \item[Portofolio program] Kumpulan seluruh berkas
          \file{.asm} yang Anda tulis, disusun menurut bab. Portofolio
          inilah bukti utama pencapaian \CPMK{03}.

    \item[Jurnal praktikum] Catatan setiap sesi praktikum yang
          memuat: tujuan, langkah, kendala yang dihadapi, cara
          mengatasinya, dan refleksi singkat. Jurnal inilah bukti
          utama pencapaian \CPMK{04} serta bahan asesmen sikap.
\end{description}

Format jurnal praktikum yang disarankan ditunjukkan pada
Tabel~\ref{tab:format-jurnal}.

\begin{table}[H]
    \centering
    \caption{Format jurnal praktikum.}
    \label{tab:format-jurnal}
    \small
    \begin{tabular}{@{}p{3.2cm}p{10.2cm}@{}}
        \toprule
        \textbf{Kolom} & \textbf{Isi} \\
        \midrule
        Sesi & Bab dan tanggal pelaksanaan \\
        \addlinespace
        Tujuan & Sub-CPMK yang hendak dicapai pada sesi ini \\
        \addlinespace
        Langkah & Urutan tindakan yang dilakukan, termasuk perintah
        perakitan \\
        \addlinespace
        Kendala & Kesalahan yang muncul beserta pesan assembler atau
        perilaku program yang salah \\
        \addlinespace
        Penyelesaian & Bagaimana kendala diatasi dan mengapa cara itu
        benar \\
        \addlinespace
        Refleksi & Apa yang dipahami lebih baik setelah sesi ini \\
        \bottomrule
    \end{tabular}
\end{table}

\subsection{Peralatan Minimum}

Anda memerlukan perangkat praktik yang dibahas pada
Bab~\ref{chap:memulai-assembly}:
TASM untuk merakit, TLINK untuk menggabungkan, DEBUG untuk memeriksa
register dan memori, emu8086 untuk menjalankan dan menelusuri program,
serta DOSBox sebagai lingkungan DOS. Turbo Debugger dipakai pada bab
debugging untuk penelusuran per instruksi.

\subsection{Konvensi Penulisan dalam Buku Ini}

\begin{itemize}
    \item Istilah asing ditulis miring pada kemunculan pertama, contoh:
          \eng{debugger}.
    \item Nama instruksi dan direktif ditulis dengan huruf tetap,
          contoh: \asm{MOV}, \dir{DB}.
    \item Nama register dan flag ditulis dengan huruf tetap, contoh:
          \reg{AX}, \flag{ZF}.
    \item Bilangan heksadesimal ditulis dengan akhiran \code{h},
          contoh: \hex{1234}; bilangan biner dengan akhiran
          \code{b}, contoh: \bin{1010}.
    \item Diagram dan tabel dirujuk dengan nomor, bukan ditulis
          manual: ``Gambar~\ref{fig:backward-design}''.
\end{itemize}
